% TOP_PROBLEM: {"id": 3325, "title": "Universal point sets for planar graphs", "category_id": 3, "category": {"id": 3, "name": "graph_theory", "display_name": "Graph Theory", "description": "Problems involving graphs, networks, and their properties.", "slug": "graph-theory", "order_index": 3, "created_at": "2026-07-31T15:26:25.671Z"}, "status": "open", "problem_number": "OPG-326", "set_id": 12}
% FIRST_POSTED: 2026-10-01
% ENTRY_KIND: statement_only
% UnsolvedMath ID 3325; source catalogue code OPG-326
% !TeX program = lualatex
% Definition and sources only; this file is not an LLM solution attempt.
\documentclass[11pt,a4paper]{article}
\usepackage[margin=25mm]{geometry}
\usepackage{fontspec}
\setmainfont{lmroman10-regular.otf}[BoldFont=lmroman10-bold.otf,
 ItalicFont=lmroman10-italic.otf,BoldItalicFont=lmroman10-bolditalic.otf]
\usepackage{amsmath,amssymb,mathtools}
\usepackage{unicode-math}
\setmathfont{latinmodern-math.otf}
\AtBeginDocument{\let\setminus\smallsetminus}
\usepackage{xurl,hyperref}
\hypersetup{colorlinks=true,urlcolor=blue}
\setcounter{secnumdepth}{0}
\setlength{\parindent}{0pt}
\setlength{\parskip}{5pt}
\setlength{\emergencystretch}{3em}
\begin{document}
\section*{Universal point sets for planar graphs}\label{3325}
{\small UnsolvedMath ID 3325 · OPG-326 · Graph Theory · OpenGarden}
\subsection{Definitions and mathematical statement}
Fix an integer $n\ge1$. A finite point set $P\subset\mathbb R^2$ is
$n$-universal for planar graphs if every finite simple planar graph
$G$ with exactly $n$ vertices has an injective map
$\varphi:V(G)\to P$ such that drawing each edge $uv$ as the straight
segment $[\varphi(u),\varphi(v)]$ gives a plane embedding. Thus edge
interiors are disjoint, no edge contains a vertex other than its
ends, and adjacent edges meet only at their common endpoint.

The graph can use any $n$ points of $P$, and the assignment of
vertices to those points is chosen separately for each graph.
The points themselves must be fixed independently of $G$.
Define $u(n)$ as the smallest possible cardinality of such a point
set. The question asks whether
\[
                  \exists C>0\ \forall n\ge1:\qquad u(n)\le Cn.
\]
Equivalently, is there a family of universal point sets whose sizes
grow at most linearly with $n$? The constant is absolute and no
particular grid, coordinate precision, or algorithm for constructing
the sets is required by the statement.

The graph comes with no prescribed embedding, so its rotation
system may be chosen along with the drawing. The straight-segment
requirement is essential: arbitrary curved edges would allow many
more placements. Allowing unused points means this is not the
stronger demand for a universal set of exactly $n$ points, and no
fixed labelling of points by graph vertices is assumed.
\subsection{Short English statement}
Is there a point set of at most Cn points that supports a crossing-free straight-line drawing of every n-vertex planar graph?
\subsection{Sources}
\begin{itemize}
\item[S1] ULAM.AI and the UnsolvedMath Contributors, supplied problems.json, record 3325 (OPG-326), OpenGarden collection. Catalogue identity and source transcription; CC BY 4.0.
\url{https://huggingface.co/datasets/ulamai/UnsolvedMath}.
\item[S2] Open Problem Garden, Universal point sets for planar graphs. Original problem and discussion; the mathematical formulation below is expanded from the supplied source record.
\url{http://www.openproblemgarden.org/op/small_universal_point_sets_for_planar_graphs}.
\end{itemize}
\end{document}
